% =============================================================================
%  Chapter 3: Methodology
%  Master Thesis – André Kuhn (MScIDS, HSLU)
%  "How well do LLMs anonymize text data?"
% =============================================================================

\chapter{Methodology}
\label{ch:methodology}

This chapter describes the experimental framework designed to answer the four
research questions. The methodology consists of three main stages:
(1)~synthetic dataset creation with ground-truth PII annotations,
(2)~anonymization using twenty in-house pipeline configurations spanning
classical and LLM-based approaches plus one external API
baseline, and (3)~multi-dimensional evaluation covering PII detection
accuracy, semantic preservation, and quality assessment.
Figure~\ref{fig:methodology_overview} illustrates the overall workflow.

\begin{figure}[htbp]
  \centering
  \begin{tikzpicture}[
    node distance=10mm and 0mm,
    every node/.style={font=\small},
    stage/.style={
      rectangle, rounded corners=10pt,
      fill=red!8, draw=red!55!black, line width=0.7pt,
      minimum width=0.55\textwidth,
      align=center, inner sep=10pt,
    },
    flow/.style={-{Stealth[length=3.5mm,width=3mm]}, line width=0.9pt,
                 draw=black!70},
  ]
    \node[stage] (gen) {\textbf{Synthetic Data Generation}};

    \node[stage, below=of gen] (anon) {%
      \textbf{Anonymization}\\[4pt]
      \begin{tabular}{@{}l@{}}
        $\bullet$ Classical tag-and-replace\\
        $\bullet$ Fine-tuned tag-and-replace\\
        $\bullet$ LLM tag-and-replace\\
        $\bullet$ Prompt-based rewriting\\
        $\bullet$ External API\\
      \end{tabular}
    };

    \node[stage, below=of anon] (eval) {%
      \textbf{Multi-Dimensional Evaluation}\\[4pt]
      \begin{tabular}{@{}l@{}}
        $\bullet$ PII Detection (P/R/F1 per tier)\\
        $\bullet$ BERTScore\\
        $\bullet$ Utility (LLM-as-Judge)\\
      \end{tabular}
    };

    \draw[flow] (gen)  -- (anon);
    \draw[flow] (anon) -- (eval);
  \end{tikzpicture}
  \caption{Overview of the methodological workflow. Synthetic customer notes
           are generated and validated, then processed by twenty in-house
           anonymization pipeline configurations (four categories) and one
           external API baseline. Results are evaluated along
           three dimensions: PII detection accuracy, semantic preservation,
           and quality assessment.}
  \label{fig:methodology_overview}
\end{figure}


% =====================================================================
\section{Synthetic Dataset Creation}
\label{sec:data_creation}
% =====================================================================

\subsection{Motivation for Synthetic Data}
\label{sec:data_motivation}

Real customer notes from financial institutions are classified as sensitive
personal data under both the FADP and the GDPR. The author's employer, a
Swiss bank, further restricts such data through an internal information
security policy that prohibits any transfer of customer records outside
its internal infrastructure, irrespective of duration or the application
of pseudonymisation. The use of authentic customer notes within a personal
computing environment was therefore not authorised. Since the experimental
design adopted in this thesis requires all anonymisation pipelines to be
executed locally on a single consumer-grade GPU, necessarily situated
outside the bank's network perimeter, access to authentic data was
precluded by the deployment constraints of the study itself. Synthetic
data generation consequently provides the only viable path for
constructing a labelled dataset with complete ground-truth PII
annotations \citep{tan2024llm_annotation}.

The use of synthetic data offers a key methodological advantage: because the
generation process controls all PII content, ground-truth annotations are
produced by design rather than through error-prone manual labeling. This enables
exact computation of precision and recall against known entity boundaries.


\subsection{Generation Process}
\label{sec:data_generation}

Figure~\ref{fig:data_pipeline} illustrates the five-stage data pipeline from
generation to the final annotated dataset.

\begin{figure}[htbp]
  \centering
  \begin{tikzpicture}[
    every node/.style={font=\footnotesize},
    box/.style={
      rectangle, rounded corners=8pt,
      draw=red!55!black, line width=0.6pt,
      minimum width=24mm, minimum height=18mm,
      align=center, inner sep=4pt, font=\footnotesize\bfseries,
    },
    sub/.style={font=\scriptsize\itshape, align=center, text width=26mm},
    flow/.style={-{Stealth[length=2.5mm,width=2mm]}, draw=black!70,
                 line width=0.7pt},
    node distance=4mm and 5mm,
  ]
    \node[box, fill=red!10]  (gen)  {Gemini 2.5 Pro\\Generation};
    \node[box, fill=red!25, right=of gen]  (det)  {Deterministic\\Validation};
    \node[box, fill=red!40, right=of det]  (sem)  {Semantic\\Validation};
    \node[box, fill=red!55, right=of sem]  (lab)  {Label Studio\\Review};
    \node[box, fill=red!50!black, right=of lab,
          text=white, draw=red!70!black]   (fin)  {Final Dataset\\2{,}542 docs};

    \draw[flow] (gen) -- (det);
    \draw[flow] (det) -- (sem);
    \draw[flow] (sem) -- (lab);
    \draw[flow] (lab) -- (fin);

    \node[sub, below=of gen] {3{,}000 docs\\3 complexity levels};
    \node[sub, below=of det] {Tag balance\\Format checks};
    \node[sub, below=of sem] {Gemini Flash\\LLM audit};
    \node[sub, below=of lab] {Manual review\\\& export};
    \node[sub, below=of fin] {50/25/25 split\\12 PII categories};
  \end{tikzpicture}
  \caption{Data generation and validation pipeline. Raw documents are generated
           by Gemini~2.5~Pro, validated in two stages, manually reviewed in
           Label~Studio, and split into train/dev/test subsets.}
  \label{fig:data_pipeline}
\end{figure}

The dataset is generated using the Gemini~2.5~Pro API \citep{google2025gemini},
which serves as the text generator. This choice ensures that the generating model
is substantially larger and more capable than the locally deployed models used
for anonymization (7--8 billion parameters), minimizing the risk that the
evaluation models have seen their own training patterns in the test data.

Each API call generates a batch of ten customer notes in German, simulating
realistic interactions between a client advisor and a corporate banking customer.
The generation prompt defines the task as producing Internal CRM Notes for a Swiss Corporate \& Institutional Clients (CIC)
relationship manager. The prompt specifies:

\begin{itemize}
    \item \textbf{Language:} Standard German (Hochdeutsch), with Swiss banking
    terminology permitted (e.g., ``Traktanden'', ``Pendenzen'').
    \item \textbf{PII annotation:} Inline XML tags for twelve entity categories
    (see Section~\ref{sec:pii_taxonomy}).
    \item \textbf{Strict annotation rules:} Nine rules governing tag syntax,
    including flat structure (no nesting), no attributes, mandatory closing tags,
    exclusion of salutations and articles from entity spans, and full inclusion
    of legal entity suffixes in organization names.
    \item \textbf{Scenario diversity:} Notes cover treasury operations, financing,
    corporate governance, client interactions, lifecycle operations, and
    administrative tasks.
    \item \textbf{Gold-standard examples:} Two fully annotated examples
    (formal and hasty style) are included in every prompt.
\end{itemize}


\subsection{Complexity Levels}
\label{sec:complexity_levels}

To evaluate robustness across varying text quality, the dataset is generated
at three complexity levels. Each level is controlled jointly by a dedicated
generation prompt, specifying register, sentence structure, abbreviations,
and (for High) intentional typos, and by the sampling temperature, with
lower temperatures producing more deterministic, formal output and higher
temperatures yielding noisier, fragmented text:

\begin{enumerate}
    \item \textbf{Low complexity} ($T = 0.30$): Formal visit reports written in
    complete sentences with perfect Standard German and an objective tone.
    These represent official records and are the easiest texts for NER systems
    to process.

    \item \textbf{Medium complexity} ($T = 0.70$): Standard CRM notes using
    concise, professional banking terminology and industry abbreviations.
    These represent the typical quality of daily advisor notes.

    \item \textbf{High complexity} ($T = 0.85$): Hasty quick-logs written in
    bullet points and sentence fragments, intentionally containing realistic
    typos (e.g., ``habne'' for ``haben'', ``fianziell'' for ``finanziell'').
    These simulate the most challenging real-world scenario where advisors
    type quickly between client calls.
\end{enumerate}

Each complexity level contributes approximately one-third of the total dataset,
enabling a per-level analysis of pipeline robustness.


\subsection{PII Taxonomy}
\label{sec:pii_taxonomy}

The annotation schema defines twelve PII categories, organized into three
tiers based on detection difficulty. The conceptual distinction between
direct identifiers (names, locations, organisations) and quasi-identifiers
(occupations, ages, nationalities, education) follows the privacy
literature originating with \citet{sweeney2002kanonymity} and adopted in
text-anonymisation benchmarks such as the Text Anonymization Benchmark
\citep{pilan2022tab} and the indirect-identifier schema of
\citet{baroud2025indirect}. This thesis extends that two-way split with
an intermediate tier for structured identifiers (IBANs, emails, phone
numbers, dates, monetary amounts), motivated by the practical observation
that such entities are reliably captured by deterministic pattern
matching, while the other two tiers require learned models. The resulting
three-tier structure is central to the evaluation framework, as it
reveals where classical methods fail and where LLMs provide clear
advantages.

\begin{table}[htbp]
\centering
\caption{PII taxonomy: twelve entity categories organized into three tiers of
detection difficulty.}
\label{tab:pii_taxonomy}
\begin{tabular}{lll}
\toprule
\textbf{Tier} & \textbf{Category} & \textbf{Description} \\
\midrule
\multirow{3}{*}{\textbf{Tier 1: Direct NER}}
  & \texttt{PER}   & Person names (excluding salutations and titles) \\
  & \texttt{LOC}   & Locations: addresses, cities, cantons, countries \\
  & \texttt{ORG}   & Organizations (including legal suffixes: AG, GmbH) \\
\midrule
\multirow{5}{*}{\textbf{Tier 2: Structured}}
  & \texttt{IBAN}  & Swiss IBANs (starting with CH) \\
  & \texttt{EMAIL} & Email addresses \\
  & \texttt{PHONE} & Phone numbers (Swiss formats) \\
  & \texttt{DATE}  & Dates (various German formats) \\
  & \texttt{MONEY} & Monetary amounts with currency \\
\midrule
\multirow{4}{*}{\textbf{Tier 3: Quasi-Identifiers}}
  & \texttt{JOB}     & Job titles and occupations \\
  & \texttt{AGE}     & Ages, birth years, age references \\
  & \texttt{NATION}  & Nationalities \\
  & \texttt{EDU}     & Educational qualifications \\
\bottomrule
\end{tabular}
\end{table}

Figure~\ref{fig:pii_tiers} visualizes the three tiers with their respective
F1-score ranges for classical and fine-tuned approaches.

\begin{figure}[htbp]
  \centering
  \begin{tikzpicture}[
    every node/.style={font=\small},
    tier/.style={
      rectangle, rounded corners=10pt,
      draw=red!55!black, line width=0.7pt,
      minimum width=42mm, minimum height=34mm,
      align=center, inner sep=6pt,
    },
    title/.style={font=\small\bfseries\itshape, text=red!55!black},
    cats/.style={font=\footnotesize, align=center},
    diffarrow/.style={-{Stealth[length=3mm,width=2.5mm]},
                       line width=0.9pt, draw=red!55!black},
    node distance=0mm and 6mm,
  ]
    % Three tier boxes with progressive shading
    \node[tier, fill=red!8]                        (t1) {};
    \node[tier, fill=red!22, right=of t1]          (t2) {};
    \node[tier, fill=red!50, right=of t2,
          text=white, draw=red!75!black]           (t3) {};

    % Box contents (placed manually so spacing matches across tiers)
    \node[title, anchor=north]
        at ($(t1.north)+(0,-3mm)$) {Tier 1 \textemdash{} Direct NER};
    \node[cats] at (t1) {PER (Person names)\\
                         LOC (Locations)\\
                         ORG (Organizations)};

    \node[title, anchor=north]
        at ($(t2.north)+(0,-3mm)$) {Tier 2 \textemdash{} Structured};
    \node[cats] at (t2) {IBAN (Bank accounts)\\
                         EMAIL (Emails)\\
                         PHONE (Phone numbers)\\
                         DATE (Dates)\\
                         MONEY (Amounts)};

    \node[title, anchor=north, text=white]
        at ($(t3.north)+(0,-3mm)$) {Tier 3 \textemdash{} Quasi-Identifiers};
    \node[cats, text=white] at (t3) {JOB (Occupations)\\
                                     AGE (Age references)\\
                                     NATION (Nationalities)\\
                                     EDU (Education)};

    % Top arrow with caption
    \draw[diffarrow] ($(t1.north west)+(2mm,7mm)$)
                  -- ($(t3.north east)+(-2mm,7mm)$)
                  node[midway, above, font=\footnotesize\itshape, text=red!55!black]
                  {Increasing detection difficulty};
  \end{tikzpicture}
  \caption{PII taxonomy organized into three tiers of increasing detection
           difficulty. Tier~1 categories are standard named entities,
           Tier~2 categories follow structured patterns, and Tier~3
           categories are contextually embedded quasi-identifiers.}
  \label{fig:pii_tiers}
\end{figure}

\textbf{Tier~1} categories correspond to standard NER labels that pre-trained
models such as spaCy and BERT are explicitly designed to detect.
\textbf{Tier~2} categories follow regular patterns (e.g., IBAN checksums, email
syntax) and can be captured by regex-based recognizers.
\textbf{Tier~3} categories are quasi-identifiers that require contextual
understanding, for example, determining that ``ETH-Ingenieur'' refers to a
job title rather than an organization. These are the categories where classical
methods are expected to fail and LLMs should demonstrate their advantage.


\subsection{Data Validation}
\label{sec:data_validation}

Generated data undergoes a two-stage validation process to ensure annotation
quality:

\textbf{Stage~1: Deterministic validation:} Rule-based checks verify the
structural integrity of annotations:
\begin{itemize}
    \item Tag balance: every opening tag has a matching closing tag.
    \item No hallucinated tags: only the twelve defined categories are accepted.
    \item No bracket leaks: entity text spans do not contain XML markers.
    \item Character offset consistency: \texttt{text[start:end]} matches the
    recorded entity value for every annotation.
\end{itemize}

\textbf{Stage~2: Semantic validation:} Gemini~Flash is used as an auditor to
identify critical tagging errors and logical inconsistencies that deterministic
rules cannot catch, such as missing entities that appear in the text but are not
annotated, or entity boundaries that include surrounding context inappropriately.

Records that fail either stage are discarded. After parse failures (9
notes) and validation rejections (449 notes, 15.0\%), 2{,}542 of the
original 3{,}000 generated notes are retained and imported into
Label~Studio for manual review, producing the definitive ground-truth
annotations used for all evaluations.


\subsection{Dataset Statistics}
\label{sec:dataset_stats}

The final annotated dataset comprises 2{,}542 customer notes. For evaluation,
the data is split into three subsets stratified by complexity level:

\begin{itemize}
    \item \textbf{Training set:} 1{,}270 documents (50\%), used for BERT and
    LLM fine-tuning.
    \item \textbf{Development set:} 634 documents (25\%), used for
    hyperparameter tuning and early stopping.
    \item \textbf{Test set:} 638 documents (25\%), used for all reported
    evaluations. Every pipeline is evaluated on this identical test set.
\end{itemize}

The stratified split ensures that each subset contains a proportional
representation of Low, Medium, and High complexity documents (test set:
218~Low, 215~Medium, 205~High), preventing distributional bias in the
evaluation. The slight imbalance across complexity levels reflects
differential rejection rates during validation: high-complexity texts with
intentional typos more frequently exhibit tag balance errors and are
therefore rejected at a higher rate during the deterministic validation
stage.


The canonical tag-and-replace ground truth for every document is produced
by replacing each annotated span with its category placeholder (e.g.,
``[PER]'', ``[MONEY]''). Representative examples of the annotated format
and the masked form, one per complexity level, are provided in
Appendix~\ref{sec:appendix_example_notes}.


% =====================================================================
\section{Anonymization Pipelines}
\label{sec:anonymization_pipelines}
% =====================================================================

A total of twenty in-house anonymization pipeline configurations are
evaluated, spanning four categories: classical baselines, a fine-tuned
BERT model, LLM-based tag-and-replace (with three prompting strategies
per model), and LLM-based prompt rewriting (with two prompting strategies
per model). In addition, the Anonymizer API, provided by the supervisor,
is included as an external baseline.
Table~\ref{tab:pipeline_overview} provides an overview. For Qwen~2.5 and SauerkrautLM, all three tag-and-replace strategies
(zero-shot, few-shot, few-shot with verification) are evaluated, but the
detailed discussion in the results chapter focuses on the best-performing
configuration per model (few-shot with verification) to maintain readability.
Full results for all configurations are provided in the appendix.

\begin{table}[htbp]
\centering
\caption{Overview of all anonymization pipeline configurations: twenty
in-house pipelines plus one external API baseline. For the LLM
tag-and-replace category, each of the three models is evaluated across
three prompting strategies (zero-shot, few-shot, few-shot+verify),
yielding nine configurations; for the LLM rewrite category, each model
is evaluated across two prompting strategies (zero-shot, few-shot),
yielding six configurations. The detailed discussion focuses on the
best configuration per model in the tag-and-replace category (marked
with $\star$).}
\label{tab:pipeline_overview}
\small
\begin{tabular}{lll}
\toprule
\textbf{Pipeline} & \textbf{Category} & \textbf{Approach} \\
\midrule
\multicolumn{3}{l}{\textit{Classical baselines}} \\
spaCy + Regex                          & Classical      & NER + pattern matching \\
BERT (pretrained) + Regex              & Classical      & NER + pattern matching \\
Presidio                               & Classical      & NER + custom recognizers \\
\midrule
\multicolumn{3}{l}{\textit{Fine-tuned encoder}} \\
BERT Fine-Tuned                        & Fine-tuned     & Token classification \\
\midrule
\multicolumn{3}{l}{\textit{LLM tag-and-replace: Llama-3}} \\
Llama-3 [zero-shot]                    & LLM tag-repl.  & Zero-shot prompting \\
Llama-3 [few-shot]                     & LLM tag-repl.  & 3-shot prompting \\
Llama-3 [few-shot+verify]$^\star$      & LLM tag-repl.  & 3-shot + self-verify \\
\midrule
\multicolumn{3}{l}{\textit{LLM tag-and-replace: Qwen2.5}} \\
Qwen2.5 [zero-shot]                    & LLM tag-repl.  & Zero-shot prompting \\
Qwen2.5 [few-shot]                     & LLM tag-repl.  & 3-shot prompting \\
Qwen2.5 [few-shot+verify]$^\star$      & LLM tag-repl.  & 3-shot + self-verify \\
\midrule
\multicolumn{3}{l}{\textit{LLM tag-and-replace: SauerkrautLM}} \\
SauerkrautLM [zero-shot]               & LLM tag-repl.  & Zero-shot prompting \\
SauerkrautLM [few-shot]                & LLM tag-repl.  & 3-shot prompting \\
SauerkrautLM [few-shot+verify]$^\star$ & LLM tag-repl.  & 3-shot + self-verify \\
\midrule
\multicolumn{3}{l}{\textit{LLM fine-tuned (QLoRA)}} \\
Llama-3 [fine-tuned]                   & LLM fine-tuned & QLoRA tag-and-replace \\
\midrule
\multicolumn{3}{l}{\textit{LLM prompt-based rewriting}} \\
Llama-3 [prompt zero-shot]             & LLM rewrite    & Zero-shot rewriting \\
Llama-3 [prompt few-shot]              & LLM rewrite    & 3-shot rewriting \\
Qwen2.5 [prompt zero-shot]             & LLM rewrite    & Zero-shot rewriting \\
Qwen2.5 [prompt few-shot]              & LLM rewrite    & 3-shot rewriting \\
SauerkrautLM [prompt zero-shot]        & LLM rewrite    & Zero-shot rewriting \\
SauerkrautLM [prompt few-shot]         & LLM rewrite    & 3-shot rewriting \\
\midrule
\multicolumn{3}{l}{\textit{External API}} \\
Anonymizer API                         & External API   & Classical NLP techniques (no LLM) \\
\bottomrule
\end{tabular}
\end{table}


% ----- Classical Baselines -----
\subsection{Classical Baselines}
\label{sec:classical_baselines}

\subsubsection{spaCy + Regex}

The first baseline combines the spaCy \texttt{de\_core\_news\_lg} model
\citep{honnibal2020spacy} for Tier~1 named entity recognition (PER, LOC, ORG)
with a set of regular expressions for Tier~2 structured entities (IBAN, EMAIL,
PHONE, DATE, MONEY). The regex patterns are designed for Swiss-specific formats,
including Swiss IBAN checksums, Swiss phone number formats (+41, 0xx), and
German date notation (DD.MM.YYYY). This pipeline cannot detect Tier~3
quasi-identifiers.

\subsubsection{BERT (Pretrained) + Regex}

The second baseline replaces spaCy's NER component with a pre-trained German
BERT model (\texttt{fhswf/bert\_de\_ner}) for token-level classification of
Tier~1 entities. The same regex patterns as the spaCy baseline handle Tier~2
entities. This pipeline likewise cannot detect Tier~3 quasi-identifiers.

\subsubsection{Microsoft Presidio}

The third baseline uses Microsoft Presidio \citep{microsoft2024presidio}, a
production-grade PII detection framework that unifies NER models with
pattern-based recognizers and context-aware confidence boosting. Presidio uses
spaCy's \texttt{de\_core\_news\_lg} model as its NLP engine for German NER.
Custom recognizers are added for Swiss-specific patterns (IBAN, phone numbers,
monetary amounts) and partial Tier~3 support through keyword-based recognizers.
A minimum confidence threshold of $0.3$ is applied to all detections.


% ----- Fine-Tuned BERT -----
\subsection{Fine-Tuned BERT}
\label{sec:bert_finetuned}

A \texttt{bert-base-german-cased} model is fine-tuned for token-level NER on
the training set using all twelve PII categories with IOB2 labeling. This is the
only classical-family model capable of detecting Tier~3 quasi-identifiers,
as it learns these categories from the annotated training data.

Training uses the following hyperparameters: 5~epochs, batch size~16,
learning rate $3 \times 10^{-5}$ with linear warmup (10\% of steps) and weight
decay of $0.01$. The best checkpoint is selected based on development set
performance. The training set comprises 1{,}270 documents (50\% of the data).


% ----- LLM Tag-and-Replace -----
\subsection{LLM-Based Tag-and-Replace}
\label{sec:llm_tag_replace}

The tag-and-replace approach, inspired by GPT-NER \citep{wang2023gptner},
instructs the LLM to copy the input text while marking entity spans with
delimiter tokens (\texttt{@@LABEL~...~\#\#}). A deterministic parser extracts
character-offset spans from the model's output. This two-step process separates
the LLM's role (entity detection) from the anonymization action (replacement
with \texttt{[LABEL]} placeholders), enabling precise evaluation against
ground-truth annotations.

Three locally deployed LLMs are evaluated:

\begin{enumerate}
    \item \textbf{Meta-Llama-3-8B-Instruct} \citep{dubey2024llama3}: A
    general-purpose instruction-tuned model (8B parameters, 16\,GB in float16).
    This serves as the primary baseline due to its strong multilingual
    capabilities.

    \item \textbf{Qwen2.5-7B-Instruct}: A 7B-parameter model from the Qwen
    family, considered one of the strongest models in its size class at the time
    of evaluation.

    \item \textbf{Llama-3.1-SauerkrautLM-8b-Instruct}: A German-specialized
    variant of Llama-3.1, fine-tuned on German instruction data by VAGOsolutions.
    This model tests whether German-specific training provides an advantage for
    German-language PII detection.
\end{enumerate}

All models run locally in float16 precision on a single NVIDIA RTX~4090 GPU
(24\,GB VRAM) without quantization.

Three prompting strategies are evaluated for each model:

\begin{itemize}
    \item \textbf{Zero-shot:} The prompt defines all twelve entity categories
    with descriptions and annotation rules, but provides no examples.

    \item \textbf{Few-shot (3 examples):} Three fully annotated German banking
    notes of varying complexity are prepended to the prompt, demonstrating the
    expected annotation format.

    \item \textbf{Few-shot + self-verification:} After the initial annotation,
    the LLM is prompted a second time to verify each extracted entity, confirming
    or rejecting it. This additional step aims to reduce false positives at the
    cost of doubled inference time.
\end{itemize}

All three strategies are evaluated for all three models, yielding nine
tag-and-replace configurations. Since few-shot with self-verification
consistently outperforms the other strategies across all models, the detailed
discussion in the results chapter focuses on this configuration. Results for
all configurations are reported in the appendix.


% ----- LLM Prompt-Based Rewriting -----
\subsection{LLM-Based Prompt Rewriting}
\label{sec:llm_prompt_rewrite}

Unlike tag-and-replace, the prompt-based rewriting approach instructs the LLM to
directly produce an anonymized version of the input text. The model is prompted
to rewrite the text such that all PII is removed or replaced with
natural-sounding alternatives, while preserving the original meaning and
structure as closely as possible.

This approach is evaluated for all three LLMs (Llama-3, Qwen~2.5, SauerkrautLM)
in a few-shot configuration. It is particularly relevant for RQ2a and RQ2b,
as it enables an assessment of semantic preservation that goes beyond simple
placeholder substitution. Because the model rewrites rather than masks, the
resulting text can be more natural and readable, but also carries a higher risk
of hallucination and semantic drift.

Since prompt-based rewrites do not produce entity-level predictions,
precision, recall, and F1 cannot be computed directly. The privacy of
these pipelines is therefore assessed through the PII leakage rate
defined in Section~\ref{sec:eval_detection}.


\subsection{External API}
\label{sec:external_api}

To contextualize the in-house pipelines against an external reference,
the Anonymizer API, provided by the supervisor, is also evaluated. The
service performs anonymization using various classical NLP techniques
rather than an LLM. Because it transmits customer data to an external
endpoint, it does not satisfy the same on-premises constraint as the
in-house pipelines and is reported as a separate category in the
results. No PII detection metrics (precision, recall, F1) are computed
for this pipeline because it returns anonymized text only; privacy is
assessed through the PII leakage rate and the LLM-judge scores.


% =====================================================================
\section{Evaluation Framework}
\label{sec:evaluation}
% =====================================================================

The evaluation framework is organized along two dimensions corresponding to
the research questions:

\begin{itemize}
    \item \textbf{RQ1, Detection and Privacy:} Did the pipeline successfully
    detect and remove PII? Assessed through entity-level precision, recall, and
    F1-score, PII leakage rates, and LLM-judge scores for anonymization quality
    and re-identification risk.

    \item \textbf{RQ2, Semantic Preservation and Utility:} Did the pipeline
    preserve everything else? Assessed through BERTScore (masked and
    full text) and LLM-judge scores for readability, meaning
    preservation, and hallucination, complemented by normalized
    Levenshtein distance, ROUGE-1, and BLEU-4 as surface-level
    indicators.
\end{itemize}


% ----- PII Detection and Privacy (RQ1) -----
\subsection{PII Detection and Privacy (RQ1a, RQ1b)}
\label{sec:eval_detection}

PII detection is evaluated at the entity level using standard NER metrics
\citep{pilan2022tab, altalla2025evaluating}. Let $\mathit{TP}$ denote the
number of true positives (predicted entities that match a ground-truth
entity in both span and category), $\mathit{FP}$ the number of false
positives (predicted entities with no corresponding ground-truth
entity), and $\mathit{FN}$ the number of false negatives (ground-truth
entities that are not predicted):

\begin{itemize}
    \item \textbf{Precision:} The proportion of predicted entities that match
    a ground-truth entity (correct span and correct category).
    \begin{equation}
        \mathrm{Precision} = \frac{\mathit{TP}}{\mathit{TP} + \mathit{FP}}
        \label{eq:precision}
    \end{equation}
    \item \textbf{Recall:} The proportion of ground-truth entities that are
    successfully detected.
    \begin{equation}
        \mathrm{Recall} = \frac{\mathit{TP}}{\mathit{TP} + \mathit{FN}}
        \label{eq:recall}
    \end{equation}
    \item \textbf{F1-score:} The harmonic mean of precision and recall.
    \begin{equation}
        F_1 = 2 \cdot \frac{\mathrm{Precision} \cdot \mathrm{Recall}}
                           {\mathrm{Precision} + \mathrm{Recall}}
        \label{eq:f1}
    \end{equation}
\end{itemize}

Entity matching follows a \emph{strict} criterion: a predicted entity is
counted as correct only if both the text span and the category label match the
ground truth exactly. Additionally, a \emph{relaxed} matching mode requires
at least 50\% character overlap of the shorter span, following the evaluation
protocol established by the i2b2/UTHealth de-identification shared task
\citep{stubbs2015i2b2}. 

Metrics are computed at four levels of granularity:
\begin{enumerate}
    \item Per individual category (all twelve labels).
    \item Per tier (Tier~1, Tier~2, Tier~3).
    \item Overall (all categories combined).
    \item Per complexity level (Low, Medium, High).
\end{enumerate}

Recall receives particular emphasis in the analysis, as anonymization systems
should minimize the omission of sensitive information, a missed PII entity
constitutes a privacy breach, whereas a false positive merely results in
unnecessary redaction \citep{pilan2022tab}.

For prompt-based rewriting pipelines, where entity-level predictions
are not produced, a \emph{PII leakage rate} serves as the privacy
metric. Inspired by \citet{siyan2025papillon}, the leakage rate
measures how many ground-truth PII strings survive in the anonymized
text:

\begin{equation}
\text{LEAK} = \frac{\text{number of PII strings found in anonymized text}}{\text{total number of PII strings}}
\label{eq:leakage}
\end{equation}

\noindent For each ground-truth entity, the metric checks whether its
text appears in the anonymized output using case-insensitive word
boundary matching. Word boundary matching is necessary for German to
avoid false positives from compound words and inflected forms; for
example, the city name ``Bern'' (LOC) must not produce a spurious
match inside ``Berner'' (the corresponding adjective, common in Swiss
banking text such as ``Berner Kantonalbank''). A leakage rate of 0 means all PII was removed; higher values indicate
more PII exposure. Nested annotations, such as a surname appearing
inside an organisation name, can produce cascading matches when the
parent entity is left unmasked.

All detection and privacy metrics are reported overall and per complexity level
(Low, Medium, High).


% ----- Semantic Preservation and Utility (RQ2) -----
\subsection{Semantic Preservation and Utility (RQ2a, RQ2b)}
\label{sec:eval_semantic}

Evaluating semantic preservation is less straightforward than detection
accuracy and requires several complementary signals. The thesis combines
two embedding-based metrics, two surface-level metrics, and the
LLM-as-judge framework described in
Section~\ref{sec:eval_judge}.

\paragraph{Masked-text BERTScore.} The primary indicator of structural
fidelity to non-PII content. PII spans in the original text are first
replaced with category placeholders (e.g.\ \texttt{[PER]}); all
placeholder tokens in both the original and the resulting anonymized
text are then masked, producing \texttt{masked\_original} and
\texttt{masked\_anonymized}. BERTScore \citep{zhang2020bertscore} is
computed between these two masked texts using a multilingual BERT model
(\texttt{bert-base-multilingual-cased}). Perfect non-PII preservation
yields a BERTScore of~1.0; any deviation indicates that the
anonymization process has altered content beyond the PII spans.

\paragraph{Full-text BERTScore.} BERTScore is additionally computed on
the unmasked original and anonymized texts. This second variant captures
overall coherence shifts that the masked variant deliberately factors
out, including the impact of placeholder tokens themselves on the
surrounding text.

\paragraph{Surface-level metrics.} Normalized Levenshtein distance on
masked texts captures small insertions or deletions that embedding-based
similarity might miss. ROUGE-1 and BLEU-4 on full texts measure unigram
overlap and 4-gram precision, respectively \citep{staab2025iclr}. These
quantify lexical fidelity but, unlike BERTScore and the judge model, do
not capture meaning equivalence under paraphrase.

All metrics in this section are reported overall and per complexity
level (Low, Medium, High). The LLM-as-judge framework, which provides
holistic readability, meaning-preservation, and hallucination scores
that complement the automated metrics above, is described next.


% ----- LLM-as-Judge -----
\subsection{LLM-as-Judge Quality Assessment}
\label{sec:eval_judge}

A capable judge model provides a holistic assessment of anonymized texts
that embedding-based similarity cannot capture in isolation. Following
the methodology of \citet{staab2025iclr}, Gemini~2.5~Flash is used as
the judge model, chosen for its ability to discriminate between
anonymization quality levels more effectively than smaller local
models.

For each anonymized document (paired with its original), the judge scores five
dimensions. These dimensions are assigned to the research questions as follows:

\paragraph{Privacy dimensions (RQ1).} These complement the automated detection
metrics by providing a holistic assessment of PII removal effectiveness:
\begin{enumerate}
    \item \textbf{Anonymization quality} (1--10): How effectively has PII been
    removed or masked?

    \item \textbf{Re-identification risk} (1--10): How difficult would it be
    to re-identify the individuals mentioned, given only the anonymized text?
\end{enumerate}

\paragraph{Utility and faithfulness dimensions (RQ2).} These assess whether
the anonymization process preserved the text's informational value:
\begin{enumerate}
    \setcounter{enumi}{2}
    \item \textbf{Readability} (1--10): How readable is the anonymized text on
    its own, without reference to the original?

    \item \textbf{Meaning preservation} (1--10): How well does the anonymized
    text preserve the original's factual and contextual meaning?

    \item \textbf{Hallucination} (0/1): Does the anonymized text contain
    information that was not present in the original? This dimension captures a
    failure mode unique to generative models that cannot be detected by
    entity-level metrics.
\end{enumerate}

Scores are aggregated per pipeline, per complexity level, and overall across all
638 test documents.


% ----- Experimental Protocol -----
\subsection{Experimental Protocol}
\label{sec:experimental_protocol}

To ensure a fair comparison, all twenty-one pipeline configurations are
evaluated on the identical held-out test set of 638~documents using fixed
random seeds. The in-house experimental pipeline, data generation,
anonymization, and evaluation, runs entirely on local hardware, in line
with the on-premises constraint motivated in
Section~\ref{sec:research_gap}. The Anonymizer API baseline is the only
component that transmits data to an external provider and is reported
separately. Concrete hardware, software versions, numerical precision,
and seed values are documented in
Section~\ref{sec:impl_experimental_setup}.
